%! Linear Regression — Unit 1, Lesson 1.5 — Slides (Beamer)
% Deck follows the lesson's gradual-release flow: targets → warm-up → hook → guided notes (I do /
% we do) → individual practice (you do alone) → debrief → close & homework. No group activity.
\documentclass[aspectratio=169, 11pt]{beamer}
\usepackage{algebra2-beamer}   % provides \forestheader{} and \sectionlabel[color]{}
% Coordinate grid: {xmin}{xmax}{ymin}{ymax}
\newcommand{\lgrid}[4]{%
  \draw[step=1, linegray!40, very thin] (#1,#3) grid (#2,#4);
  \draw[->, thick, charcoal] (#1,0)--(#2,0) node[below right, font=\tiny]{$x$};
  \draw[->, thick, charcoal] (0,#3)--(0,#4) node[left, font=\tiny]{$y$};}
\newcommand{\dpt}[2]{\fill[navy] (#1,#2) circle (2.8pt);}
% The free-throw data: hours practised (x) vs. free throws made out of 15 (y).
\newcommand{\ftdata}{\dpt{1}{4}\dpt{2}{4}\dpt{3}{7}\dpt{4}{7}\dpt{5}{11}\dpt{6}{11}\dpt{7}{12}\dpt{8}{14}}

\begin{document}

% TITLE SLIDE (CourseName is not defined in beamer, so the name is literal)
{
\setbeamercolor{background canvas}{bg=forest}
\begin{frame}[plain]
  \vspace{2.8cm}\hspace{0.55cm}%
  \begin{minipage}{0.9\textwidth}
    {\color{goldacc}\bfseries\small LESSON 1.5}\\[0.5cm]
    {\color{white}\bfseries\LARGE Linear Regression}\\[1.2cm]
    {\color{forestmid}\small Algebra 2~$\cdot$~Unit 1: Linear Functions}
  \end{minipage}
\end{frame}
}

% LEARNING TARGETS
\begin{frame}[t, plain]
  \forestheader{Today's targets}
  \sectionlabel{What you will be able to do}
  \begin{itemize}\setlength{\itemsep}{6pt}
    \item Read a \textbf{scatterplot} and describe its association: \emph{direction},
          \emph{form}, \emph{strength}.
    \item Interpret the \textbf{correlation coefficient} $r$ --- sign for direction, size for
          strength --- and match an $r$ to a plot.
    \item Read a \textbf{line of best fit} $\hat y=ax+b$, interpret its slope and intercept
          \emph{in context}, and \textbf{predict} with it.
    \item Tell \textbf{interpolation} from \textbf{extrapolation}, judge a prediction, and say
          why \textbf{correlation is not causation}.
  \end{itemize}
  \vspace{4pt}
  \begin{block}{How today runs}
    Warm-up (5) $\to$ notes together (20) $\to$ \textbf{practice on your own} (15) $\to$ debrief
    (10) $\to$ close \& start the homework (10).
  \end{block}
\end{frame}

% WARM-UP
\begin{frame}[t, plain]
  \forestheader{Warm-Up: three things you already know}
  \sectionlabel{5 minutes --- alone, then check with a neighbour}
  \begin{enumerate}\setlength{\itemsep}{7pt}
    \item \textbf{Read the points.} Coordinates of $A$, $B$, $C$, $D$. Do they lie on one line?
          Do they \emph{mostly} rise?
    \item \textbf{Predict with a rule.} $y=4x+10$ at $x=2$, $5$, $8$.
    \item \textbf{Slope and direction.} For $y=3x-5$: slope, intercept, and what happens to $y$
          each time $x$ goes up by $1$.
  \end{enumerate}
  \vspace{6pt}
  \begin{block}{Hold on to this}
    No single line goes through all four points --- but they \emph{mostly} rise. That is the
    difference between a \emph{rule} and a \emph{trend}. Keep it.
  \end{block}
\end{frame}

% HOOK
\begin{frame}[t, plain]
  \forestheader{Hook: hours practised, free throws made}
  \sectionlabel[goldacc]{Eight players, one week, 15 shots each}
  \begin{columns}[T]
    \begin{column}{0.40\textwidth}
      \centering
      \begin{tikzpicture}[scale=0.26]
        \lgrid{-0.5}{9}{-0.5}{15.5}
        \ftdata
      \end{tikzpicture}
    \end{column}
    \begin{column}{0.56\textwidth}
      \begin{itemize}\setlength{\itemsep}{5pt}
        \item As hours go up, what do the free throws do? \; \textbf{They tend to go up.}
        \item Could a single line pass through every point? \; \textbf{No.} \; Could one line
              summarise the trend? \; \textbf{Yes.}
        \item $5$ hours made $11$; $6$ hours made $11$. Does more practice \emph{guarantee}
              more baskets? \; \textbf{No --- it \emph{tends} to.}
      \end{itemize}
    \end{column}
  \end{columns}
  \vspace{4pt}
  \begin{block}{The idea that carries today}
    No exact rule fits messy data --- but a \textbf{line of best fit} can summarise the trend.
    Today we read it, measure how well it fits, and predict with it.
  \end{block}
\end{frame}

% NOTES 1 — ASSOCIATION
\begin{frame}[t, plain]
  \forestheader{Notes 1: scatterplots \& association}
  \sectionlabel{Direction, form, strength --- no numbers yet}
  \begin{center}
  \begin{tikzpicture}[scale=0.30]
    \lgrid{-0.5}{7}{-0.5}{7}
    \dpt{1}{2}\dpt{2}{2}\dpt{3}{4}\dpt{4}{4}\dpt{5}{6}\dpt{6}{6}
    \node[font=\scriptsize\bfseries] at (3,-1.8){(a)};
  \end{tikzpicture}
  \hspace{0.6cm}
  \begin{tikzpicture}[scale=0.30]
    \lgrid{-0.5}{7}{-0.5}{7}
    \dpt{1}{6}\dpt{2}{5}\dpt{3}{5}\dpt{4}{3}\dpt{5}{2}\dpt{6}{2}
    \node[font=\scriptsize\bfseries] at (3,-1.8){(b)};
  \end{tikzpicture}
  \hspace{0.6cm}
  \begin{tikzpicture}[scale=0.30]
    \lgrid{-0.5}{7}{-0.5}{7}
    \dpt{1}{3}\dpt{2}{6}\dpt{3}{2}\dpt{4}{5}\dpt{5}{3}\dpt{6}{5}
    \node[font=\scriptsize\bfseries] at (3,-1.8){(c)};
  \end{tikzpicture}
  \end{center}
  \begin{itemize}\setlength{\itemsep}{4pt}
    \item \textbf{Direction:} as $x$ rises, does $y$ rise (positive), fall (negative), or
          neither? \; (a) positive, (b) negative, (c) none.
    \item \textbf{Form:} do the points follow a \emph{line}? \; (a) and (b) yes; (c) no.
    \item \textbf{Strength:} how tightly? \; (a) and (b) strong.
  \end{itemize}
  \vspace{2pt}
  The free-throw data: positive, linear, strong.
\end{frame}

% NOTES 2 — r  (the target misconception)
\begin{frame}[t, plain]
  \forestheader{Notes 2: the correlation coefficient $r$}
  \sectionlabel[redacc]{Sign is direction; size is strength}
  \begin{center}
  \begin{tikzpicture}[scale=1.1]
    \draw[thick, charcoal, ->] (-3.6,0)--(3.6,0);
    \foreach \x/\lab in {-3/-1, -1.5/-0.5, 0/0, 1.5/0.5, 3/1}{\draw (\x,0.12)--(\x,-0.12) node[below, font=\scriptsize]{$\lab$};}
    \node[font=\scriptsize, forest] at (-3,0.45){strong negative};
    \node[font=\scriptsize, forest] at (0,0.45){no linear trend};
    \node[font=\scriptsize, forest] at (3,0.45){strong positive};
  \end{tikzpicture}
  \end{center}
  \begin{columns}[T]
    \begin{column}{0.48\textwidth}
      \textbf{$r=0.9$}: positive, strong --- points hug a line that \emph{rises}.
    \end{column}
    \begin{column}{0.48\textwidth}
      \textbf{$r=-0.9$}: negative, strong --- points hug a line that \emph{falls}.
    \end{column}
  \end{columns}
  \vspace{4pt}
  Match: $0.90\to$ (a), \; $-0.85\to$ (b), \; $0.05\to$ (c).
  \vspace{4pt}
  \begin{block}{Careful --- which is the stronger fit, $r=-0.9$ or $r=0.4$?}
    $r=-0.9$: \; $|{-0.9}|=0.9>0.4$. The minus sign says which way the trend runs; the $0.9$
    says how tightly. Negative is \emph{not} weak.
  \end{block}
\end{frame}

% NOTES 3 — THE LINE
\begin{frame}[t, plain]
  \forestheader{Notes 3: the line of best fit --- and what its numbers mean}
  \sectionlabel{Technology gives $\hat y=1.5x+2$, \; $r=0.97$}
  \begin{columns}[T]
    \begin{column}{0.40\textwidth}
      \centering
      \begin{tikzpicture}[scale=0.26]
        \lgrid{-0.5}{9}{-0.5}{15.5}
        \draw[very thick, forest] (0,2)--(8.6,14.9);
        \ftdata
      \end{tikzpicture}
    \end{column}
    \begin{column}{0.56\textwidth}
      \begin{itemize}\setlength{\itemsep}{5pt}
        \item \textbf{Slope $1.5$:} each additional \emph{hour} of practice is associated with
              about $1.5$ more free throws --- a rate, with units.
        \item \textbf{Intercept $2$:} the predicted free throws at $0$ hours --- a baseline.
        \item \textbf{$r=0.97$:} strong, because $|r|$ is close to $1$.
        \item The hat on $\hat y$: a \emph{predicted} value, not a data point.
      \end{itemize}
    \end{column}
  \end{columns}
  \vspace{4pt}
  The line runs through the \emph{middle} of the points, not through them --- technology picks the
  line whose predictions are closest to the data overall.
\end{frame}

% NOTES 4 — PREDICTING AND ITS LIMITS
\begin{frame}[t, plain]
  \forestheader{Notes 4: predicting --- and its limits}
  \sectionlabel{And your turn}
  \begin{columns}[T]
    \begin{column}{0.48\textwidth}
      \textbf{Inside the data: $x=6$}
      \[ \hat y = 1.5(6)+2 = 11 \]
      \textbf{Interpolation.} Reasonable.
    \end{column}
    \begin{column}{0.48\textwidth}
      \textbf{Outside the data: $x=12$}
      \[ \hat y = 1.5(12)+2 = 20 \]
      \textbf{Extrapolation.} Only $15$ shots --- impossible.
    \end{column}
  \end{columns}
  \vspace{3pt}
  \begin{block}{Two limits}
    \textbf{The line does not know the story:} outside the data it keeps going; reality does not.
    \textbf{Correlation is not causation:} ice-cream sales and drownings both rise in summer
    ($r>0$) --- hot weather, a \emph{lurking variable}, drives both.
  \end{block}
  \vspace{3pt}
  \textbf{Guided Practice, together:} exercise minutes vs.\ resting heart rate, $\hat y=-0.4x+78$,
  $r=-0.88$ --- negative \emph{and} strong; the slope with units; $\hat y(30)=66$; interpolation;
  and no, a strong $r$ does not prove cause.
\end{frame}

% INDIVIDUAL PRACTICE LAUNCH
\begin{frame}[t, plain]
  \forestheader{Individual Practice: on your own}
  \sectionlabel[goldacc]{15 minutes $\cdot$ silent $\cdot$ last page of your notes}
  \begin{enumerate}\setlength{\itemsep}{6pt}
    \item \textbf{Match $r$ to the plot.} $0.95$, $0.10$, $-0.70$ --- and which fit is strongest?
    \item \textbf{Interpret the model.} Experience vs.\ salary (in thousands):
          $\hat y=2.5x+38$. Slope with units, intercept, $\hat y(10)$, and interpolation or
          extrapolation.
    \item \textbf{Where the line breaks.} Hot cocoa vs.\ temperature, $\hat y=-3x+260$ over
          $20$--$60^\circ$. Predict at $40^\circ$ and at $90^\circ$ --- then say what went wrong.
  \end{enumerate}
  \vspace{5pt}
  \begin{block}{The rule}
    Direction from the sign, strength from the size; every slope gets units; check every
    prediction against the data's range. Stuck? Look \emph{up the page} at the notes --- not at
    your neighbour. Done early? Add hot days to problem~3 --- what happens to $r$?
  \end{block}
\end{frame}

% DEBRIEF
\begin{frame}[t, plain]
  \forestheader{Debrief: correct your own page in a second colour}
  \sectionlabel[redacc]{10 minutes}
  \begin{enumerate}\setlength{\itemsep}{5pt}
    \item \textbf{1 --- read $r$ in two steps.} Only $R$ falls, so $R$ is $-0.70$. $P$ hugs its
          line: $0.95$. $Q$ has no line to hug: $0.10$. Strongest: $P$.
    \item \textbf{The must-land moment.} $r=-0.9$ vs.\ $r=0.4$: which points sit closer to
          their line? The minus sign is direction; the $0.9$ is strength. Negative is not weak.
    \item \textbf{2 --- units, or it is not an interpretation.} \$2{,}500 more salary per
          additional year. $\hat y(10)=63$ is \$63{,}000 --- interpolation. $x=30$ is twice the
          longest career measured --- extrapolation.
    \item \textbf{3 --- the line does not know the story.} $\hat y(90)=-10$ cups. The line hits
          zero near $87^\circ$, far outside $20$--$60^\circ$. Conclude ``very hot days sell almost
          none'' --- the model's job ends where the data ends.
  \end{enumerate}
\end{frame}

% CLOSE
\begin{frame}[t, plain]
  \forestheader{Close}
  \sectionlabel{What changed today}
  You walked in with functions whose rules were \emph{exact}. You walk out able to read a
  \emph{trend} in scattered data, measure it with $r$, interpret the line of best fit in context,
  predict with it --- and know where to stop trusting it.
  \vspace{5pt}
  \begin{block}{Homework --- start it now}
    The \textbf{Homework} page in your packet. We start item~1 together and you keep going; whatever
    is left is due next class. Six problems plus an extension: matching $r$ to plots, a
    sign-versus-size pair, a model from a data table, the special values of $r$, a causation
    critique, and one multiple-choice item.
  \end{block}
  \vspace{4pt}
  \textbf{Next:} the \textbf{Unit 1 Test} --- then Unit 2, where the data can \emph{curve} and the
  curve of best fit becomes a parabola.
\end{frame}

\end{document}
